% 자동 생성 — scripts/build_tables.py. 손으로 고치지 않는다.
% 출처: output/results/18_cancellations.json : 이중차분.2025-03-24 재지정
\begin{table}[htbp]
  \centering
  \caption{취소거래가 규제 직후에 몰리는가 (2025-03-24 재지정).}
  \label{tab:cancel}
  \begin{tabular}{lrrrrr}
    \toprule
    \textbf{창 (개월)} & \textbf{처치 전} & \textbf{처치 후} & \textbf{통제 전} & \textbf{통제 후} & \textbf{DID (\%p)} \\
    \midrule
    3 & 4.38 & 8.18 & 7.04 & 11.25 & -0.4 \\
    6 & 4.05 & 8.96 & 6.81 & 10.02 & +1.7 \\
    12 & 3.40 & 7.03 & 5.20 & 7.91 & +0.9 \\
    \bottomrule
  \end{tabular}

  \begin{minipage}{\linewidth}\footnotesize
  \textit{주.} 창을 바꾸면 부호가 바뀐다. 처치군 특유의 집중이 없다. 취소를 다시 넣어도 주 결과는 $0.3\%$p 움직인다.
  \end{minipage}
\end{table}
